\documentclass[10pt]{article}
\usepackage[margin=20mm]{geometry}
\usepackage{amsmath,amssymb,fontspec,unicode-math,xurl,hyperref}
\IfFontExistsTF{DejaVu Serif}{\setmainfont{DejaVu Serif}}{\setmainfont{Latin Modern Roman}}
\setmathfont{latinmodern-math.otf}
\setlength{\parindent}{0pt}\setlength{\parskip}{6pt}\sloppy
\title{Technical Verification of CUPS Coverage Guarantees}
\author{AI-assisted technical review}
\date{}
\begin{document}
\maketitle
Abstract: CUPS improves several reported pose-estimation metrics, but three mathematical steps used for its uncertainty guarantees fail under their stated assumptions. Label-dependent calibration weights can destroy coverage even for independent identically distributed observations. Ranking these weights does not give the claimed temporal-decay bound. A separate Beta-distance bound is numerically and analytically false. These are major defects in core theoretical derivations; they do not establish that the reported empirical improvements are false. Scientific adjudication remains pending.\par
\section*{Identity and inspected scope}
Harry Zhang and Luca Carlone, CUPS: Improving Human Pose-Shape Estimators with Conformalized Deep Uncertainty, ICML 2025, I25-43692 (0323), OpenReview `v7ebdKPTnG`. The frozen publisher PDF has 19 pages; PDF and printed page numbering agree. The main text and Appendices A–H were read. Detailed verification concerns Definitions 1–3, the weighted quantile, Theorems 1–3 and their proofs, and the interpretation of Table 3. Original architecture, pose and coverage figures on pp. 4, 6–8 and the decisive proof pages 14 and 16 were viewed. The acquired packet contains no author repository; a separate OpenReview attachment inventory was inaccessible. No author program, pose estimator, dataset or training run was executed.\par
The relevant cited conformal paper was consulted at the authors’ public preprint, specifically Sections 3.1 and 5.5. Its fixed-weight requirement and its conditional-total-variation alternative matter here. Other cited literature was not independently audited. No later CUPS erratum comparison was performed.\par
\section*{Method and checks that survive}
CUPS uses a learned nonconformity score \(S(X,Y)=\sigma(\phi(X,Y))\), a sigmoid applied to an MLP embedding, and a set \(\{Y:S(X,Y)\le\tau\}\). The adversarial losses push scores of ground-truth poses toward zero and those of predictions toward one; the predictor receives the opposing score loss. This is consistent with using lower scores as greater conformity. Training does not impose a deterministic mathematical constraint on the score’s distribution.\par
The independent arithmetic check of Table 1 agrees with the stated 3DPW improvements over GLoT: PA-MPJPE improves by \(1.9\) mm, MPJPE by \(4.5\) mm, and MPVPE by \(4.6\) mm. Acceleration error increases from \(6.6\) to \(6.9\), as the text acknowledges. These checks concern printed numbers, not benchmark reproduction. Appendix F’s finite proposal aggregation is a practical heuristic; no guarantee of sampling uniformly from the entire conformal set is inferred.\par
\section*{Substantive findings}
Problem 1 (Theorem 1 does not cover the proposed label-dependent weights).  This is a major defect in the core coverage justification. Theorem 1 (p. 6 and Appendix B, pp. 12–13), combined with Definition 3 (pp. 6, 14), uses an unconditional swap-distance correction for weights computed from each calibration response.\par
Equation reference: eq:weight. Definition 3 specifies \[
 w_i=\exp(-\|\phi(X_i,\widehat Y_i)-\phi(X_i,Y_i)\|^2/T).\] These depend on the ground-truth calibration label \(Y_i\). The borrowed fixed-weight theorem does not automatically apply. This is more than a missing citation condition: the coverage assertion can fail under exchangeability, where every unconditional swap distance is zero.\par
To see this, fix \(X\) and a predictor \(\widehat Y=0\), and let a scalar pose coordinate \(U\) be uniform on \([0,1]\), independently across the calibration and test records. Keep other coordinates fixed. Choose a fixed MLP projection \(\phi(X,Y)=U\), so \(S=\sigma(U)\) and \(\phi(X,\widehat Y)=0\). A linear projection is within the specified MLP score family. Set \(T=0.01\) and \(\alpha=0.1\). This realizes [eq:weight] (\#eq:weight) exactly with \(w_i=e^{-U_i^2/T}\in(0,1]\). It is a counterexample to the distribution-free mathematical claim, not an assertion about a particular trained CUPS checkpoint.\par
Because \(\sigma\) is strictly increasing, the weighted \(0.9\) quantile of scores is the sigmoid of the weighted \(0.9\) quantile of \(U\). The law of large numbers gives its limit \(q\) uniquely by \[\int_0^q e^{-u^2/T}\,du
 =0.9\int_0^1e^{-u^2/T}\,du,
 \qquad
 q=\sqrt T\,\operatorname{erf}^{-1}
       (0.9\operatorname{erf}(1/\sqrt T))=0.1163087154.\] Strict positivity of the density gives quantile convergence. Conditional test coverage equals the realized threshold in \(U\) coordinates; bounded convergence therefore gives limiting marginal coverage \(q\), far below \(0.9\). The sample and its coordinate-swapped version have identical joint laws, so the unconditional total-variation terms in Theorem 1 are all zero. Consequently the claimed bound cannot hold for all finite calibration sizes under these assumptions.\par
Equation (5) omits the extra atom at \(+\infty\) that appears in the Appendix B proof and the cited method. Including an atom of mass \(1/(1+\sum_iw_i)\) repairs that finite-sample convention, but does not repair this example: the atom vanishes as \(n\to\infty\), leaving the same limit. An original diagnostic calculation with 500 independent calibration sets of size 1500 gives mean conditional coverage \(0.11618\) without this atom and \(0.11855\) with it. The decisive argument is the limit above, not the Monte Carlo estimate.\par
The cited conformal analysis explicitly fixes weights before observing the calibration records. It also describes a different bound using total variation conditional on random weights, averaged over those weights; this does not reduce to zero merely because the unconditional records are exchangeable. Treating training as fixed does not fix weights that depend on new calibration labels. A defensible repair is to specify weights in advance and include the test atom, or to prove and control the appropriate conditional distances. Neither repair is established for the current feature-distance procedure. This invalidates the advertised calibration guarantee without refuting the empirical effectiveness of the adversarial score.\par
Problem 2 (Feature-weight rank is substituted for chronological order).  Theorem 2, equation (11), and its Appendix C.1 proof (p. 14) assert a \(\rho^k\) bound after a changepoint, using \(w'_i=\rho^{n+1-\pi(w_i)}\), where \(\pi(w_i)\) is the ranked position of feature weight \(w_i\), rather than time index \(i\).\par
Take \(n=4,k=3,\rho=1/2\). Let the sole old observation have the largest feature weight, with ranks \((4,1,2,3)\). Then \[(w'_1,w'_2,w'_3,w'_4)=(1/2,1/16,1/8,1/4).\] Let the old score be deterministically \(0.2\), and the recent and test scores be deterministically \(0.8\). Recent full observations can be i.i.d., with continuous feature weights below the old weight; their ordering is immaterial because the sum of geometric rank weights is fixed. For example the old weight can be \(0.9\) and recent weights lie in \((0.1,0.3)\), all realizable in Definition 3 by squared embedding distances \(-\log w_i\) at \(T=1\). Fixed ground-truth embeddings \(\operatorname{logit}(0.2)\) and \(\operatorname{logit}(0.8)\) give these scores; predicted embeddings can differ by \(\sqrt{-\log w_i}\). The recent records and the test record may have identical distributions, as the changepoint assumption requires.\par
Swapping a recent score with the test score leaves the score vector unchanged, so its total variation is zero. Swapping the old score with the test score changes a deterministic vector, so its total variation is one. Thus the asserted left side is \[\frac{1/2}{1/2+1/16+1/8+1/4}=\frac8{15}
 >\frac18=\rho^k.\] Even if the proof’s denominator \(1+\sum_iw'_i\) replaces the theorem’s \(\sum_iw'_i\), it equals \(8/31>1/8\). Ties among recent ranks are avoided by continuous weights. The construction satisfies even the stronger condition that recent score-vector swaps have distance zero; it does not rely on confusing marginal and joint exchangeability.\par
Equation (19)’s geometric-series step silently needs the large weights to belong to recent time indices. No assumption relates feature rank to time this way. With genuine recency weights \(w'_i=\rho^{n+1-i}\), the normalized old mass is instead \[\frac{\sum_{i\le n-k}\rho^{n+1-i}}
 {\sum_{i\le n}\rho^{n+1-i}}
 =\rho^k\frac{1-\rho^{n-k}}{1-\rho^n}\le\rho^k,
 \qquad 0<\rho<1.\] This verifies only the repaired geometric bound, conditional on the post-change swap-distance assumptions. It does not verify coverage for the original label-dependent method. The defect is Major because it invalidates the special-purpose bound used to justify Table 3.\par
Problem 3 (The Beta-distance bound is false near a parameter boundary).  Theorem 3, equations (12) and (28), and Appendix C (pp. 15–17) assert a bound depending only on \(n\) and \(|a_1-a_2|\le k\) for \(\mathrm{Beta}(a_1,n-a_1)\) and \(\mathrm{Beta}(a_2,n-a_2)\).\par
Even granting the charitable scalar interpretation of the notation \(S_\theta(Z)\), take \(n=100,k=1,a_1=1,a_2=2\). These are positive integer Beta parameters satisfying every stated parameter constraint. For \(P=\mathrm{Beta}(1,99)\) and \(Q=\mathrm{Beta}(2,98)\) the density ratio is \(f_Q/f_P=98x/(1-x)\), so there is exactly one crossing at \(x_*=1/99\). Integration gives \[d_{\rm TV}(P,Q)=F_P(x_*)-F_Q(x_*)
 =\left(\frac{98}{99}\right)^{98}=0.3697484374.\] The paper’s bound is \[\sqrt{2-2(1-2/101)^{1/2}}=0.1410708778,\] strictly smaller. This refutes the pairwise inequality from which the weighted result is derived; no assumption on weights can make that pairwise inequality valid. The notation also alternates between a score vector and a univariate Beta variable, but resolving that ambiguity cannot fix the false scalar inequality.\par
The decisive error is visible in (27). The quadratic \(D(a_1)=a_1n-a_1^2-a_1k+kn\) is maximized at \((n-k)/2\). The function \(1-kn/D\) increases in positive \(D\). Substituting this maximum therefore gives an upper bound, whereas the proof uses it as a lower bound. In fact with \(a_1=1,a_2=2,k=1\) and \(n\to\infty\), the exact variation tends to \(e^{-1}\), while the claimed bound tends to zero. A restriction that keeps Beta parameters away from the boundary would need a new proof and is absent here.\par
Equation (22)’s exact Hellinger affinity remains correct. A valid, parameter-dependent replacement is \[d_{\rm TV}(P,Q)\le\min\left\{1,
 \sqrt{2-2\frac{B((a_1+a_2)/2,n-(a_1+a_2)/2)}
 {\sqrt{B(a_1,n-a_1)B(a_2,n-a_2)}}}\right\}.\] This follows by Cauchy–Schwarz applied to \(|p-q|=|\sqrt p-\sqrt q|(\sqrt p+\sqrt q)\); it is \(0.4793803\) in the example, consistent with the exact distance. It does not recover the published \(k\)-only guarantee. This is Major for the core theoretical uncertainty bound.\par
\section*{Impact, limitations, and disposition}
Table 3 reports weighted empirical coverage \(86.2\%\), \(87.3\%\) and \(89.0\%\) for nominal \(90\%\). These observed quantities are separate from the disputed claimed lower bounds. The counterexamples do not change the recorded percentages or establish a benchmark fabrication. They establish that the present derivations cannot certify the proposed coverage behavior. Labels/activity changes used to estimate \(k\) in Remark 2 also do not by themselves establish a bound on fitted Beta parameters; no empirical identification of that link was reproduced.\par
Dimension | State\par
Supported overall assessment | Major\par
P1 / P2 / P3 | Major / Major / Major\par
Verified limited repairs | Fixed-weight scope; recency sum; exact affinity\par
Benchmark reproduction | Not performed\par
Scientific adjudication | Pending\par
Human review / author adoption | None recorded / unknown\par
2 H. Zhang and L. Carlone, CUPS: Improving Human Pose-Shape Estimators with Conformalized Deep Uncertainty, ICML 2025. Frozen publisher PDF, 19 pp. https://openreview.net/forum?id=v7ebdKPTnG. R. Foygel Barber, E.J. Candès, A. Ramdas and R.J. Tibshirani, Conformal Prediction Beyond Exchangeability, Annals of Statistics 51 (2023), 816–845. Consulted authors’ 62-page preprint, Sections 3.1 and 5.5. https://www.stat.cmu.edu/\textasciitilde{}ryantibs/papers/nexcp.pdf.\par
\section*{Scope and limitations}
The accompanying independent Python script computes the Beta CDF difference, Hellinger affinity, rank-weight masses, weighted-quantile limit and seeded diagnostic. It uses Python 3.13.9, NumPy 2.4.3 and SciPy 1.16.3; no author source was imported. Hash-bound page images, coverage, findings, compilation and visual-inspection records accompany the report. The subsequent focused independent acceptance and corrected source locations are recorded in the existing verification files; no human review is recorded.\par
\end{document}
